\chapter*{Rappel - Les fonctions}

\section*{Repérer les fonctions}

\underline{Exercice 1 :}\\

Voici des extraits du code de la bibliothèque Python « openai » (openai-python 0.x, © OpenAI, licence MIT), repérer les définitions et les appels de fonctions (vous pouvez surligner en différentes couleurs).
Dans chacun des cas, sans chercher à comprendre la fonction, déterminer les paramètres et les valeurs de retours.\\
\begin{lstlisting}[language=Python]
   
def _console_log_level():
    if openai.log in ["debug", "info"]:
        return openai.log
    elif OPENAI_LOG in ["debug", "info"]:
        return OPENAI_LOG
    else:
        return None

def logfmt(props):
    def fmt(key, val):
        # Handle case where val is a bytes or bytesarray
        if hasattr(val, "decode"):
            val = val.decode("utf-8")
        # Check if val is already a string to avoid re-encoding into ascii.
        if not isinstance(val, str):
            val = str(val)
        if re.search(r"\s", val):
            val = repr(val)
        # key should already be a string
        if re.search(r"\s", key):
            key = repr(key)
        return "{key}={val}".format(key=key, val=val)

    return " ".join([fmt(key, val) for key, val in sorted(props.items())])


def get_object_classes():
    # This is here to avoid a circular dependency
    from openai.object_classes import OBJECT_CLASSES

    return OBJECT_CLASSES


def convert_to_openai_object(
    resp,
    api_key=None,
    api_version=None,
    organization=None,
    engine=None,
    plain_old_data=False,
):
    # If we get a OpenAIResponse, we'll want to return a OpenAIObject.

    response_ms: Optional[int] = None
    if isinstance(resp, openai.openai_response.OpenAIResponse):
        organization = resp.organization
        response_ms = resp.response_ms
        resp = resp.data

    if plain_old_data:
        return resp
    elif isinstance(resp, list):
        return [
            convert_to_openai_object(
                i, api_key, api_version, organization, engine=engine
            )
            for i in resp
        ]
    elif isinstance(resp, dict) and not isinstance(
        resp, openai.openai_object.OpenAIObject
    ):
        resp = resp.copy()
        klass_name = resp.get("object")
        if isinstance(klass_name, str):
            klass = get_object_classes().get(
                klass_name, openai.openai_object.OpenAIObject
            )
        else:
            klass = openai.openai_object.OpenAIObject

        return klass.construct_from(
            resp,
            api_key=api_key,
            api_version=api_version,
            organization=organization,
            response_ms=response_ms,
            engine=engine,
        )
    else:
        return resp

def merge_dicts(x, y):
    z = x.copy()
    z.update(y)
    return z
\end{lstlisting}

\section*{Trouver l'erreur}

\underline{Exercice 2 :}\\

Trouvez l'erreur dans la définition de la fonction suivante, expliquez pourquoi cela cause une erreur et formulez une correction.

\begin{lstlisting}[language=Python]
    def bonjour()
        print("Bonjour")
\end{lstlisting}

\underline{Exercice 3 :}\\

Trouvez l'erreur dans la définition de la fonction suivante, expliquez pourquoi cela cause une erreur et formulez une correction.

\begin{lstlisting}[language=Python]
    def decompte10():
    print("Debut du decompte")
        for i in range(10):
            print(10-i)
        print(0, "!")
        print("Fin du decompte")
\end{lstlisting}

\underline{Exercice 4 :}\\

Trouvez l'erreur dans la définition de la fonction suivante, expliquez pourquoi cela cause une erreur et formulez une correction.

\begin{lstlisting}[language=Python]
    def decompte(n):
        if n > 0:
        print("Vous devez donner un nombre positif")
        print("Debut du decompte")
        for i in range(n):
            print(10-i)
        print(0, "!")
        print("Fin du decompte")
\end{lstlisting}


\section*{Compréhension d'algorithmes}

\underline{Exercice 5 :}\\
Que fait la fonction suivante ? dérouler l'exécution du code ligne par ligne dans un tableau.
\begin{lstlisting}[language=Python]
    def mystere(n):
        a = 1
        for i in range(1, n+1):
            a = a*i
        return a

    mystere(5)
\end{lstlisting}

\underline{Exercice 6 :}\\
Que fait la fonction suivante ? dérouler l'exécution du code ligne par ligne dans un tableau.
\begin{lstlisting}[language=Python]
    def mystere(a, b):
        if a > b :
            temp = a
            a = b
            b = temp
        return a, b

    print(mystere(46, 12))
\end{lstlisting}

\underline{Exercice 7 :}\\
Que fait la fonction suivante ? dérouler l'exécution du code ligne par ligne dans un tableau.
\begin{lstlisting}[language=Python]
    def mystere(tab, n):
        for i in range(len(tab)) :
            tab[i] = tab[i] * n
        return tab

    print(mystere([3, 6, 10, 1, 0], 5))
\end{lstlisting}

\underline{Exercice 8 :}\\
Que fait la fonction suivante ? dérouler l'exécution du code ligne par ligne dans un tableau.
\begin{lstlisting}[language=Python]
    def mystere(tab):
        for i in range(1, len(tab)+1) :
            if tab[i-1] != tab[-i]:
                return False
        return True

    print(mystere([3, 6, 10, 1, 0, 4, 10, 6, 3]))
\end{lstlisting}

\underline{Exercice 9 :}\\
Que fait la fonction suivante ? dérouler l'exécution du code ligne par ligne dans un tableau.
\begin{lstlisting}[language=Python]
    from math import sqrt # import de la fonction racine carree
    
    def mystere(point1, point2):
        return sqrt((point1[0] - point2[0])**2 + (point1[1] - point2[1])**2)

    def mystere2(tab, depart):
        point = tab[0]
        min_dist = mystere(point, depart)
        for i in range (1, len(tab)):
            if mystere(tab[i], depart) < min_dist:
                point = tab[i]
                min_dist = mystere(tab[i], depart)
        return point

    print(mystere2([[7, 9], [2, 5], [5, 2]], [0, 0]))
\end{lstlisting}

\newpage
\section*{Définition de fonction}

\subsection*{Fonction simple}

Pour chacun des exemples ci-dessous, vous imaginerez un script de test permettant de mettre à l'épreuve votre fonction et déroulerez votre algorithme sur l'exemple de test.

\underline{Exercice 10 :}\\
Écrire une fonction \textit{alphabet} écrivant les lettres de l'alphabet dans l'ordre en minuscule dans le terminal.

\underline{Exercice 11 :}\\
Écrire une fonction \textit{factorielle10} qui fait le produit des nombres entre 1 et 10 ($1 \times 2 \times ... \times 9 \times 10)$ et affiche le résultat dans le terminal.

\subsection*{Avec paramètres}

\underline{Exercice 12 :}\\
Écrire une fonction \textit{factorielle} qui fait le produit des nombres entre 1 et n ($1 \times 2 \times ... \times n-1 \times n)$, $n$ étant un entier passé en paramètres de la fonction.

\underline{Exercice 13 :}\\
Écrire une fonction \textit{nombre\_de} qui compte le nombre d'occurrences d'une valeur $n$ donnée en paramètre dans un tableau $tab$ également donnée en paramètres et l'affiche dans le terminal.

\subsection*{Avec valeur de retour}

\underline{Exercice 14 :}\\
Écrire une fonction \textit{puissance} prenant deux paramètres $x$ et $n$ et qui calcule le nombre $x^n$ et le renvoi dans le programme principal.

\underline{Exercice 15 :}\\
Écrire une fonction qui reçoit une chaîne de caractères en paramètre et qui retourne le nombre de voyelles de cette chaine.

\underline{Exercice 16 :}\\
Écrire une fonction qui résout une équation du second degré $ax^2+bx+c=0$ avec les paramètres $a$, $b$ et $c$ entrés en paramètres et qui retourne les solutions s'il y en a.

\begin{tcolorbox}[enhanced,
    colback=green!25!black!10!white,colframe=green!75!black,title=Rappel,
    drop fuzzy shadow,watermark color=white]
Il faut d'abord calculer le déterminant $\Delta = b^2-4ac$ et en fonction de sa valeur, retourner les solutions correspondantes.
\begin{itemize}
    \item Si $\Delta<0$ il n'y a pas de solution.
    \item Si $\Delta>0$ il y a deux solutions $\frac{-b+\sqrt{\Delta}}{2a}$ et $\frac{-b-\sqrt{\Delta}}{2a}$.
    \item Si $\Delta=0$ il n'y a qu'une solution $\frac{-b}{2a}$.
\end{itemize} 
\end{tcolorbox}




